\scanpage{15}
\mspage{(13)}

e) Le complexe de chaines $C_\bullet(\mathcal X_*,\deleted{\mathbb Z}\inserted{k})$ permet de reconstruire tous les complexes de cochaines $C^\bullet(\mathcal X_*,k')$, pour $k'$ \inserted{$k$-Alg.} variable, par
\[
C^\bullet(\mathcal X_*,k)\simeq\operatorname{Hom}_{\mathbb Z}(C_\bullet(\mathcal X_*,\deleted{\mathbb Z}\inserted{k}),k).
\]
\inserted{(L'objet ``le plus fin'' est le $C_\bullet(\mathcal X_*,\mathbb Z)$, on a donc $C_\bullet(\mathcal X_*,k)\simeq C_\bullet(\mathcal X_*,\mathbb Z)\otimes_{\mathbb Z}k$)}
Il est possible de même de définir une \unsure{cobigèbre} \inserted{diff} $C_{\bullet\bullet}^{\mathrm{DRpd}}(\mathcal X_*,\deleted{\mathbb Z}\inserted{k})$ \inserted{$k$-coaugmentée} à copuiss. div., telle que l'on ait, pour $\forall$ $k$-Alg. $k'$ :
\[
C^{\bullet\bullet}_{\mathrm{DRpd}}(\mathcal X_*,k')\simeq\operatorname{Hom}_k(C_{\bullet\bullet}^{\mathrm{DRpd}}(\mathcal X_*,k),k')
\]
(compatible avec \unsure{ttes} les structures).

On aura d'ailleurs
\[
C_{\bullet\bullet}^{\mathrm{DRpd}}(\mathcal X_*,k')\simeq C_{\bullet\bullet}^{\mathrm{DRpd}}(\mathcal X_*,k)\otimes_k k'.
\]
L'objet le plus fin est $C_{\bullet\bullet}^{\mathrm{DRpd}}(\mathcal X_*,\mathbb Z)$. C'est lui qu'il convient de considérer (au lieu de son ``dual'' $C^{\bullet\bullet}_{\mathrm{DRpd}}(\mathcal X_*,\mathbb Z)$) si on veut aborder b) c) d) sans condts de finitude.

Notons que [\unsure{$\forall$ comme|tt comme} $C^{\unsure{*|\bullet\bullet}}_{\mathrm{DRpd}}(\mathcal X_*,k)$, pour $\mathcal X_*$ variable, transforme $\varinjlim$ quelconq. en $\varprojlim$] $C_{\bullet\bullet}^{\mathrm{DRpd}}(\mathcal X_*,\mathbb Z)$ \deleted{transforme} \inserted{commute aux} $\varinjlim$ quelconques (pas néc. \unsure{filtrantes}).

\mspage{14}
